\documentclass[11pt]{article}

\usepackage[a4paper,margin=24mm]{geometry}
\usepackage{fontspec}
\usepackage{xcolor}
\usepackage[hidelinks]{hyperref}

\setmainfont[
  Path=../latex-kiee-review-2022/fonts/,
  Extension=.otf,
  UprightFont=*-Regular,
  BoldFont=*-Bold
]{NotoSerifCJKkr}

\definecolor{reviewercolor}{RGB}{38,82,126}
\definecolor{responsecolor}{RGB}{36,113,78}
\definecolor{revisioncolor}{RGB}{145,82,24}

\setlength{\parindent}{0pt}
\setlength{\parskip}{0.55em}
\setlength{\fboxsep}{7pt}
\setlength{\emergencystretch}{3em}
\renewcommand{\familydefault}{\rmdefault}

\newcommand{\fieldheading}[2]{%
  \par\medskip\noindent
  {\color{#1}\rule{3pt}{1.25em}}\hspace{0.55em}\textbf{#2}\par\smallskip}

\newenvironment{reviewitem}[1]
  {\subsection*{#1}}
  {\par\medskip\hrule\bigskip}
\newenvironment{reviewercomment}
  {\fieldheading{reviewercolor}{1. 리뷰어 코멘트}%
   \begin{quote}}
  {\end{quote}}
\newenvironment{authorresponse}
  {\fieldheading{responsecolor}{2. 저자 답변}}
  {\par}
\newenvironment{manuscriptrevision}[1]
  {\fieldheading{revisioncolor}{3. 논문 반영 결과}%
   \textbf{수정 위치:} #1\par\smallskip}
  {\par}

\newcommand{\manuscripttitle}{비전-언어 자율주행 모델의 실행 가드 기반 궤적 후보 선택: 안전 제약 인과 연쇄}
\newcommand{\manuscriptid}{2026-D-D04-0036}
\newcommand{\revisionround}{1차 수정}
\newcommand{\responsedate}{2026년 9월 4일}

\begin{document}

\begin{center}
  {\LARGE\bfseries 심사 의견 답변서\par}
  \vspace{1em}
  {\large\bfseries \manuscripttitle\par}
\end{center}

\vspace{0.8em}
\begin{tabular}{@{}ll@{}}
논문 접수번호: & \manuscriptid \\
수정 차수: & \revisionround \\
답변일: & \responsedate
\end{tabular}

\section*{답변에 앞서}

편집위원장님과 리뷰어 여러분께,

원고를 세심하게 검토하고 건설적인 의견을 주셔서 감사합니다. 저자들은 모든 의견을
검토하여 원고를 수정했습니다. 아래에는 각 리뷰어 코멘트, 이에 대한 저자 답변,
그리고 원고에 실제로 반영한 내용과 수정 위치를 순서대로 제시합니다. 쪽 번호는
개정 원고를 기준으로 합니다.

이번 개정에서는 후보 정확도와 계약 쌍 정확도의 관계를 명확히 하고, 제약의 의미와
입력 위치의 효과를 분리하여 설명했습니다. 또한 고난도 사례와 결정적 단위 정규화
분석을 추가하고 실행 차단장치의 개입 비용을 보고했습니다. 기존 주행 비전-언어
모델 및 실행시간 검증 연구와의 차이, 국문 요약과 영문 초록의 구체성, 원고의 논리적
흐름, 학습된 가드 준수와 실제 실행 안전의 구분도 보강했습니다. 모델 외부 검증과
평가 지표의 상호 보완적 역할도 명시했습니다. 리뷰어 1에 따른 수정은 파란색,
리뷰어 2에 따른 수정은 보라색으로 표시했습니다.

\section*{리뷰어 1}

\begin{reviewitem}{R1-G1}
  \begin{reviewercomment}
    본 논문은 기존 CoC에 실행 제약을 통합한 Safety-Constrained CoC와
    paired-contract 평가 방법을 제안한 연구로, 주제의 참신성과 학문적 의미가
    인정되며 전반적인 논문 구성도 체계적인 편이다. 제안 방법은 표준 평가에서
    우수한 성능을 보였으나, 미관측 수치 조건에서의 일반화 성능 저하와 일부
    실험에서 영상 정보의 기여가 명확하지 않은 점은 추가적인 검증이 필요하다.
    따라서 제안 방법의 유효성을 보다 명확히 확인할 수 있도록 관련 실험과
    분석을 보완한 후 게재 여부를 다시 검토하는 것이 적절하다고 판단된다.

    본 논문은 기존 CoC가 행동의 목적과 이유는 설명하지만 행동의 허용 범위를
    명시적으로 표현하기 어렵다는 점에 착안하여, execution guard를 CoC 내부에
    통합한 Safety-Constrained CoC를 제안한다. 특히 동일한 장면과 행동 목표에서
    guard만 변경하는 paired-contract 평가를 통해 모델이 실제 제약조건에 따라
    궤적 후보를 변경하는지를 검증하였다. 제안 방법은 표준 평가에서 높은 guard
    compliance를 보였으나, 미관측 시간값 및 일부 표현 변화에서는 성능 저하가
    관찰되었다.
  \end{reviewercomment}

  \begin{authorresponse}
    지적하신 두 사항을 구분하여 검증하고 결과의 의미를 다시 정리했습니다.

    \textbf{첫째, 미관측 수치에 대한 일반화입니다.} 시간 과제에서는 학습과 표준
    평가에 사용한 상충 구역 해제 시각 2.0, 3.0, 4.0초와 겹치지 않도록 미관측
    평가에 2.5, 3.5, 4.5초를 사용했습니다. 또한 학습하지 않은 해제 안정시간과
    후보 진입 지연을 사용하여, 모델이 본 숫자를 반복하는지 시간의 대소 관계를
    적용하는지를 확인했습니다. 제약 통합 CoC는 표준 평가에서 후보 정확도와 계약
    쌍 정확도 1.000, 가드 위반률 0.000을 기록했지만, 미관측 시간값에서는 후보
    정확도 0.635, 계약 쌍 정확도 0.306, 가드 위반률 0.351로 낮아졌습니다. 이는
    새로운 시간값에 대한 일반화가 확인되었다는 결과가 아니라, 표준 분포에서 배운
    대응이 일반적인 시간 비교 규칙으로 충분히 확장되지 않았음을 보여주는 부정적
    결과입니다. 따라서 표준 평가의 높은 성능을 수치 추론 능력으로 확대하지 않고
    현재 방법의 한계로 명시했습니다.

    \textbf{둘째, 영상 정보의 기여입니다.} 저장된 세 난수 초기값의 동일 어댑터와
    동일 평가 문제를 유지하고, 영상만 원본ㆍ정보가 없는 단색 영상ㆍ시간 또는 행동
    정보가 텍스트와 충돌하는 영상으로 바꾸어 재평가했습니다. 표준 시간 과제의 후보
    정확도/계약 쌍 정확도는 원본 1.000/1.000에서 단색 영상 0.670/0.403, 충돌 영상
    0.608/0.340으로 낮아졌습니다. 따라서 이 과제에서는 합성 연속 장면 영상에 표시된
    시간 정보가 실제 선택에 사용되었습니다. 반면 정지ㆍ차선 변경의 표준 평가와
    미관측 수치 평가는 단색 영상과 반대 행동 영상에서도 두 정확도가 모두 1.000을
    유지했습니다. 이 과제에서는 행동 종류와 후보 물리량이 문장에 직접 제시되어
    영상 없이도 정답을 선택할 수 있었던 것입니다. 이에 따라 영상의 기여를 모든
    과제에 공통된 장점으로 주장하지 않고, 표준 시간 과제에서만 확인된 과제 의존적
    결과로 제한했습니다. 또한 이 영상이 글자가 포함된 합성 자료이므로 실제 도로
    영상의 시각적 근거화를 입증한 결과가 아님을 명시했습니다.
  \end{authorresponse}

  \begin{manuscriptrevision}{제5절(7--11쪽), 표 9ㆍ11과 그림 5, 제6.5--6.6절(14쪽), 제7절(15쪽)}
    \begin{enumerate}
      \item 제5절의 ``평가 이름을 읽는 법''에서 표준 평가, 동일 의미 문장 변형
      평가, 미관측 수치 평가를 구분했습니다. 특히 미관측 수치 평가는 문장 틀은
      유지하되 학습에 사용하지 않은 속도ㆍ거리ㆍ시간값을 사용하는 평가임을
      명시했습니다.

      \item 제5.4절과 표 9 및 그림 5(a)에 표준 시간 평가와 미관측 시간값 평가를
      나란히 제시했습니다. 표준 평가의 후보 정확도/계약 쌍 정확도/가드 위반률
      1.000/1.000/0.000이 미관측 시간값에서 0.635/0.306/0.351로 변한 결과를
      추가하고, 이를 ``일반적인 시간 비교 규칙으로 충분히 일반화되지 않음''으로
      해석했습니다. 따라서 미관측 수치 결과를 성공 사례가 아니라 RQ4에 대한
      부정적 답변으로 기록했습니다.

      \item 제5.4--5.5절과 표 11에 저장 어댑터의 영상 절제 재평가를 반영했습니다.
      모델과 평가 문장은 고정하고 영상만 원본ㆍ단색ㆍ충돌 조건으로 바꾼 절차를
      설명했습니다. 표준 시간 과제에서는 원본 1.000/1.000에서 단색
      0.670/0.403과 충돌 0.608/0.340으로 낮아졌지만, 정지ㆍ차선 변경의 표준 및
      미관측 수치 평가에서는 세 영상 조건 모두 1.000/1.000을 유지한 결과를 함께
      제시했습니다. 이를 통해 전자는 합성 영상의 시간 정보에 의존하고 후자는 주로
      문장 정보로 해결된다는 차이를 명시했습니다.

      \item 제6.5절 ``주 실험이 실제로 확인한 것'', 제6.6절 ``결과를 해석할 때의
      실험적 한계''와 제7절 결론에서 주장의 범위를 수정했습니다. 현재 결과는 합성
      표준 분포에서의 제약 기반 후보 선택과 과제별 영상 사용을 뒷받침하지만,
      새로운 수치 관계의 일반적 해석이나 실제 도로 영상의 시각적 근거화 및 실제
      차량 안전을 입증하지 않는다고 명시했습니다.
    \end{enumerate}
  \end{manuscriptrevision}
\end{reviewitem}

\subsection*{주요 의견}

\begin{reviewitem}{R1-M1}
  \begin{reviewercomment}
    Paired-contract accuracy의 실제 의미를 조금 더 명확하게 설명할 필요가 있다.
    Paired-contract 평가는 본 논문의 중요한 장점이지만, 두 계약을 모두 맞혀야
    성공으로 판단하는 지표이므로 일반적인 candidate accuracy보다 난도가 높다.
    따라서 두 지표의 관계와 paired accuracy가 추가적으로 측정하는 능력을
    직관적인 예시를 통해 설명하면 제안 평가법의 독창성이 더욱 잘 드러날 것으로
    생각된다.
  \end{reviewercomment}

  \begin{authorresponse}
    좋은 지적에 감사드립니다. 개정 원고에서는 두 지표의 평가 단위를 명확히
    구분하였습니다. Candidate accuracy는 $N$개 장면 쌍에서 만들어진 $2N$개 개별
    계약 문제의 정답률이고, paired-contract accuracy는 한 장면의 완화ㆍ엄격
    계약을 모두 맞힌 쌍의 비율입니다. 예를 들어 두 장면 쌍에서 세 개의 개별
    계약만 맞히면 candidate accuracy는 $3/4=0.75$이지만 paired-contract
    accuracy는 $1/2=0.50$입니다. 특히 두 계약에 같은 후보를 반복하는 모델은 각
    쌍의 한 문제를 맞혀 candidate accuracy 0.50을 얻을 수 있지만
    paired-contract accuracy는 0입니다. 따라서 후자는 단순히 더 엄격한 정확도가
    아니라, 장면ㆍ목표ㆍ후보가 동일할 때 guard 변화에 따라 선택을 올바르게
    전환하는 능력을 측정합니다.
  \end{authorresponse}

  \begin{manuscriptrevision}{제4.5절, 5쪽, 모델 외부 궤적 검증과 평가 지표}
    Candidate accuracy와 paired-contract accuracy의 분모와 성공 조건을
    명시하고, 2개 장면 쌍의 $0.75$ 대 $0.50$ 예 및 같은 후보 반복 시 $0.50$ 대
    $0$이 되는 예를 추가했습니다. 해당 개정은 파란색으로 표시했습니다.
  \end{manuscriptrevision}
\end{reviewitem}

\begin{reviewitem}{R1-M2}
  \begin{reviewercomment}
    Constraint의 ‘의미’를 학습한 효과와 ‘제시 위치/형식’의 효과를 보다 명확하게
    분리할 필요가 있다. Constraint-integrated CoC가 separate natural-language
    constraint보다 높은 성능을 보이지만, 두 조건은 constraint의 내용뿐 아니라
    입력 내 위치와 구성 방식도 다르다. 논문에서도 이러한 가능성을 일부 인정하고
    있으므로, 동일한 문장과 학습량을 유지하면서 constraint의 위치만 변경하는
    추가 ablation이 있으면 제안 방식의 효과를 보다 직접적으로 설명할 수 있으리라
    사료된다.
  \end{reviewercomment}

  \begin{authorresponse}
    의견에 동의합니다. 기존 실험의 Separate natural-language constraint 조건은
    바로 이 목적의 position-only ablation이었지만, 원고에서 통제 변수를 충분히
    설명하지 못했습니다. 두 조건은 동일 guard 문자열, 장면, 후보, 정답, 384개
    학습 예, 3 epoch 및 최적화 설정을 사용하고, guard를 기본 CoC 뒤에 이어
    붙이는지 별도 필드에 두는지만 다릅니다. 표준 시간 평가에서 통합/별도 조건의
    candidate accuracy는 1.000/0.698, paired-contract accuracy는 1.000/0.396,
    violation rate는 0.000/0.191이었습니다. 이는 현재 시간 과제에서 위치ㆍ구획
    방식이 큰 영향을 주었음을 보여줍니다. 의미 대응은 위치가 같은 별도 조건과
    shuffled 조건을 비교하여 분리했으며, paired-contract accuracy는
    0.396/0.021이었습니다. 다만 시간 과제에 inline-shuffled 조건이 없어 완전한
    $2\times2$ 요인 설계는 아니므로, 통합 조건의 차이를 의미 학습만의 효과로
    해석하지 않도록 주장을 제한했습니다. 정적 과제의 통합/별도 paired accuracy는
    1.000/0.986으로 위치 효과가 작았다는 과제 의존성도 함께 보고했습니다.
  \end{authorresponse}

  \begin{manuscriptrevision}{제4.2절(4쪽), 제5.4절(8쪽), 제6.2절(11쪽)}
    입력 조건 표에 position-only ablation임을 명시하고, 동일하게 유지한 문장,
    데이터 수, 학습량과 최적화 설정을 기술했습니다. 위치 효과와 의미 대응 효과를
    별도의 정량 대비로 제시하고, 완전한 요인 설계가 아니라는 한계와 과제 의존성을
    논의에 추가했습니다. 해당 개정은 파란색으로 표시했습니다.
  \end{manuscriptrevision}
\end{reviewitem}

\begin{reviewitem}{R1-M3}
  \begin{reviewercomment}
    표준 평가에서 반복적으로 1.000의 성능이 나타나는 점은 오히려 benchmark
    난이도에 대한 추가 검토가 필요해 보인다. 제안 방법은 여러 표준 평가에서
    candidate accuracy와 paired-contract accuracy가 1.000에 도달한다. 이는 방법의
    효과를 보여주는 동시에 현재 표준 평가가 모델 간 차이를 충분히 구분하기에는
    쉬울 가능성도 의미한다. Guard 간 차이가 작거나 여러 제약조건이 동시에
    활성화되는 hard-case를 추가하면 방법의 강점을 더욱 설득력있게 풀어낼 수 있을
    것으로 보인다.
  \end{reviewercomment}

  \begin{authorresponse}
    지적하신 대로 반복적인 1.000은 표준 평가가 모델 간 차이를 구분하기에 쉬울
    가능성을 의미합니다. 개정 원고에서는 기존 복합 표현 변화 평가를 hard-case로
    명시하고 구성을 구체화했습니다. 정지 hard-case는 정지선, 감속도, 가가속도
    세 제약을 동시에 활성화하며, 완화 계약의 두 수치 상한을 빠른 후보의
    측정값과 같게 두어 판정 여유가 0인 equality boundary를 포함합니다. 차선
    변경도 최소 간격을 빠른 후보와 같은 값으로 두고 간격 미달 시 중단ㆍ재시도
    조건을 함께 적용합니다. 세 seed의 행동별 재집계에서 제안 조건의
    candidate/paired accuracy는 정지 0.979/0.958, 차선 변경 0.917/0.833으로,
    표준 평가의 1.000보다 낮았습니다. 따라서 표준 평가는 in-distribution 학습
    확인으로 한정하고, hard-case가 ceiling을 드러낸다는 해석을 추가했습니다.
    다만 수치 간격과 표현 변화가 완전히 분리된 설계는 아니므로 이를 후속
    benchmark 과제로 명시했습니다.
  \end{authorresponse}

  \begin{manuscriptrevision}{제5.5절(9쪽), 표 9(10쪽), 제6.6절(13쪽)}
    동시 활성 제약, equality boundary와 평가 규모를 hard-case 정의에 추가하고,
    정지와 차선 변경의 행동별 성능 행을 표에 추가했습니다. 표준 1.000을 난이도
    충분성의 증거로 사용하지 않도록 논의와 한계를 수정했습니다. 해당 개정은
    파란색으로 표시했습니다.
  \end{manuscriptrevision}
\end{reviewitem}

\begin{reviewitem}{R1-M4}
  \begin{reviewercomment}
    Closed-loop 결과에서는 safety뿐 아니라 shield intervention cost도 함께 분석할
    필요가 있다. Contract+shield가 unseen risk에서 violation을 크게 감소시키지만,
    이미 학습된 상황에서는 위반 감소 없이 평균 2.539회의 개입이 발생한다. 따라서
    단순 violation/collision 감소뿐 아니라 불필요한 intervention 또는 주행
    효율성과의 trade-off를 함께 평가하면 실제 시스템 관점에서 더욱 의미 있는
    결과가 될 것이라 예상된다.
  \end{reviewercomment}

  \begin{authorresponse}
    의견에 동의하여 safety benefit과 intervention cost를 함께 재분석했습니다.
    동일 episode bank에서 shield가 없는 대응 계약과 비교하여 episode당 개입
    횟수, 평균 진행거리 변화와 평균 절대 가가속도 변화를 추가했습니다.
    Phase-complete ID/OOD에서는 양 조건 모두 위반과 충돌이 0이었지만 shield가
    각각 1.269/2.539회 개입했고, 진행거리는 0.076/0.078 m 감소하며 평균 절대
    가가속도는 0.065/0.040 증가했습니다. 반면 release-reversal에서는 9.935회
    개입과 함께 동적 계약 대비 violation rate가 0.166, collision rate가 0.125
    감소했고, 진행거리 0.030 m 감소와 평균 절대 가가속도 0.306 증가가
    동반되었습니다. 모든 completion rate는 1.000이었습니다. 따라서 unseen
    reappearance에서는 안전 이득과 비용이 함께 있었지만, 이미 이진 안전지표가
    포화된 bank에서는 추가 안전 이득 없이 비용만 발생했다고 해석했습니다.
    개별 개입의 필요성을 판정하는 독립 oracle은 없으므로 intervention precision을
    계산했다고 주장하지 않고 이를 후속 과제로 남겼습니다.
  \end{authorresponse}

  \begin{manuscriptrevision}{제5.6절(10--11쪽), 표 11(12쪽), 제6.4절(12쪽), 제7절(14쪽)}
    Intervention, progress 및 jerk 변화량을 안전 이득과 함께 보고하는 표를
    추가했습니다. 안전 이득이 없는 phase-complete 조건과 이득이 있는
    release-reversal 조건을 분리하여 해석하고, 비용 대리 지표와 독립 개입
    oracle의 부재를 명시했습니다. 논의와 결론에도 safety--efficiency trade-off를
    반영했으며 해당 개정은 파란색으로 표시했습니다.
  \end{manuscriptrevision}
\end{reviewitem}

\subsection*{세부 의견}

\begin{reviewitem}{R1-m1}
  \begin{reviewercomment}
    Static, temporal, stopping/lane-change 실험에서 서로 다른 평가 지표와 설정이
    사용되므로, 각 실험의 목적과 핵심 비교 대상을 하나의 간단한 표로 정리하면
    전체 실험 구조를 이해하는 데 도움이 될 것으로 사료된다.
  \end{reviewercomment}

  \begin{authorresponse}
    좋은 제안에 감사드립니다. 실험마다 모델, 입력 변화 및 지표가 달라 독자가
    수치를 하나의 성능 순위처럼 읽을 가능성이 있었습니다. 이를 방지하기 위해
    R1 사전 진단, 최소 학습, 정적 후보, 시간 조건, 정지/차선 변경 및 폐루프 보조
    실험을 한 표에 정리했습니다. 각 행에는 연구 질문과 평가 설정, 핵심 비교,
    주요 지표 및 해당 실험이 제공하는 증거의 범위를 함께 제시했습니다. 또한
    행 사이의 절대 수치는 비교하지 않고 각 행의 핵심 비교 안에서 해석해야 함을
    명시했습니다.
  \end{authorresponse}

  \begin{manuscriptrevision}{제5절, 표 6(8쪽)}
    실험 전체의 읽기 안내 표를 연구 질문ㆍ평가 설정, 핵심 비교, 주요 지표 및
    증거 범위 중심으로 재구성했습니다. Static, temporal, stopping/lane-change의
    상이한 설정을 구분하고 closed-loop 실험이 별도 소형 네트워크의 보조 증거임도
    표시했습니다. 해당 개정은 파란색으로 표시했습니다.
  \end{manuscriptrevision}
\end{reviewitem}

\begin{reviewitem}{R1-m2}
  \begin{reviewercomment}
    복합 표현 변화에서 발생한 오류가 대부분 cm 표현과 함께 나타난 점은 흥미로운
    결과이다. 단순 limitation으로 처리하기보다 unit normalization을 적용했을 때
    성능이 얼마나 회복되는지 추가 분석하면 논문의 practical implication을 강화할
    수 있다고 보인다.
  \end{reviewercomment}

  \begin{authorresponse}
    제안에 따라 저장된 세 seed의 어댑터를 동일한 288개 복합 표현 변화 문제에서
    재평가했습니다. 이때 72개 cm 입력의 수치와 단위만 동등한 m 값으로 결정적으로
    변환하고 장면, 후보, 정답, 나머지 문장 및 모델 가중치는 유지했습니다. cm
    부분집합의 정답 수는 57/72에서 60/72로 증가했으며, 회복된 3건은 모두 기존의
    안전하지 않은 정지 선택이었습니다. 전체 후보 정확도는 0.948에서 0.958,
    paired-contract accuracy는 0.896에서 0.917로 증가하고 guard violation rate는
    0.010에서 0.000으로 낮아졌습니다. 회귀한 정답과 비-cm 입력의 예측 변화는
    없었습니다. 다만 seed 42의 안전하지만 과도하게 보수적인 차선 변경 오답
    12건은 남았습니다. 따라서 unit normalization은 수치 안전 오류를 제거했지만
    복합 문장의 보수적 선택까지 해결하지는 못했다고 제한적으로 해석했습니다.
  \end{authorresponse}

  \begin{manuscriptrevision}{제5.5절과 표 9(9--10쪽), 제6.3절(12--13쪽)}
    단위 정규화 재평가의 통제 조건과 전후 성능을 표 9 및 결과 문단에
    추가하고, 회복된 위반 3건과 남은 보수적 오류 12건을 분리했습니다. 논의에는
    결정적 단위 정규화가 모델 외부 검증 전의 유용한 전처리이지만 일반적인 언어
    강건성 해결책은 아니라는 practical implication을 추가했습니다. 해당 개정은
    파란색으로 표시했습니다.
  \end{manuscriptrevision}
\end{reviewitem}

\begin{reviewitem}{R1-m3}
  \begin{reviewercomment}
    Dynamic contract와 monotone-stage contract의 성능 차이는 흥미롭지만 상대적으로
    설명이 짧다. 위험 재등장 상황에서 contract state의 reversible update가 왜
    중요한지 조금 더 강조하면 closed-loop 실험의 의미가 명확해질 것이라 판단된다.
  \end{reviewercomment}

  \begin{authorresponse}
    의견에 동의합니다. 해제 후 재위험 상황의 핵심은
    $\mathrm{HOLD}\rightarrow\mathrm{RELEASED}\rightarrow\mathrm{HOLD}$ 전이입니다.
    출발 허용 뒤 위험이 다시 나타나면 이전 허용을 취소하고 대기 가드를
    재활성화해야 합니다. Dynamic contract는 현재 위험에 따라 이 역전이를
    표현하지만, monotone-stage contract는 한 번 도달한 RELEASED 상태에 고정됩니다.
    동일 release-reversal bank에서 dynamic contract는 monotone-stage contract보다
    violation rate를 16.5\%p(37.9\%에서 21.4\%), collision rate를
    11.6\%p(25.9\%에서 14.3\%) 낮췄습니다. 두 조건의 목표 완료율은 모두
    100\%였습니다. 또한 계약이 0.4초 늦으면 dynamic contract의 위반률이
    48.1\%로 악화되므로 상태의 가역성뿐 아니라 최신성도 필요함을 강조했습니다.
    다만 이는 별도 소형 폐루프 환경의 메커니즘 증거임을 유지했습니다.
  \end{authorresponse}

  \begin{manuscriptrevision}{제5.6절(10--11쪽), 제6.4절(13쪽)}
    Closed-loop 비교 방법에 가역 상태전이를 명시하고, dynamic/monotone 조건의
    위반률 및 충돌률 차이를 해당 상태 갱신 메커니즘과 연결했습니다. 논의에는
    출발 허용을 영구 단계가 아닌 현재 관측에 의존하는 임시 허가로 해석해야 하며,
    위험 재등장 시 RELEASED에서 HOLD로 복귀해야 한다는 설명을 추가했습니다.
    해당 개정은 파란색으로 표시했습니다.
  \end{manuscriptrevision}
\end{reviewitem}

\clearpage
\section*{리뷰어 2}

\begin{reviewitem}{R2-C1}
  \begin{reviewercomment}
    기존 비전-언어 자율주행 모델과 본 연구의 차별성 및 기존 연구의 문제점에 대한
    설명이 부족하며 연구동기가 불분명함. 즉 본연구의 필요성 및 기존 연구의
    문제점에 대한 설명이 부족하여 연구 동기가 불분명함.
  \end{reviewercomment}

  \begin{authorresponse}
    지적에 감사드립니다. 개정 원고에서는 기존 연구를 두 흐름으로 나누어 연구
    공백을 명확히 했습니다. Driving VLM/CoC는 장면의 원인, 행동 이유 및 목표를
    설명하지만 동일 목표 후보의 허용 경계만 바꾸었을 때 선택이 전환되는지를
    직접 평가하지 않습니다. 반면 형식 검증기와 runtime monitor는 주어진 수치
    명세를 정확히 검사하지만 자연어 가드를 VLM의 학습 입력으로 사용하고 그
    의미 반영 여부를 측정하지 않습니다. 본 연구는 기존 계획기나 검증기를
    대체한다고 주장하지 않습니다. CoC에 자연어 실행 가드를 결합하고, 장면ㆍ목표ㆍ
    후보를 고정한 paired-contract와 모델 외부 수치 검증을 연결하는 중간
    인터페이스를 제안한다는 차별점을 명시했습니다.
  \end{authorresponse}

  \begin{manuscriptrevision}{제1절(1--2쪽), 제2.4절과 표 2(3--4쪽)}
    서론에 두 연구 흐름 사이의 공백과 필요한 인터페이스를 설명하는 동기 문단을
    추가했습니다. 관련 연구에는 driving VLM/CoC, formal monitor/shield 및 본
    연구의 주요 역할, 실행 조건, 진단 평가와 증거 범위를 비교하는 표를
    추가했습니다. 해당 개정은 보라색으로 표시했습니다.
  \end{manuscriptrevision}
\end{reviewitem}

\begin{reviewitem}{R2-C2}
  \begin{reviewercomment}
    전체적인 구성에 있어 논리적 연결성이 미비함. 전체적으로 논리적인 흐름을
    정리할 필요가 있음.
  \end{reviewercomment}

  \begin{authorresponse}
    의견에 따라 원고의 논리적 연결을 ``연구 공백--문제 정의--입력 및 평가 방법--
    단계별 실험--증거 범위''의 순서로 명시했습니다. 서론 말미에 절별 roadmap을
    추가하고, 문제 정의에서는 행동 목표와 가드 만족을 분리한 이유를 먼저
    설명했습니다. 방법 절은 입력, 대조군과 paired-contract, 모델 외부 판정의 세
    질문으로 구성했으며, 실험 절은 사전 진단에서 정적ㆍ시간ㆍ다중 주행 및 별도
    폐루프 분석으로 범위를 넓히는 순서를 명시했습니다. 각 실험은 확인 목적,
    비교 방법, 결과 및 의미의 동일한 구조로 읽히도록 점검했습니다.
  \end{authorresponse}

  \begin{manuscriptrevision}{제1절(2쪽), 제3절(3쪽), 제4절(4쪽), 제5절(7쪽)}
    절별 roadmap과 문제 정의ㆍ방법ㆍ실험 절의 전환 문단을 추가했습니다. 이
    문단들은 앞 절의 질문이 다음 절의 수식, 평가 절차 및 실험 순서로 어떻게
    이어지는지 명시합니다. 해당 개정은 보라색으로 표시했습니다.
  \end{manuscriptrevision}
\end{reviewitem}

\begin{reviewitem}{R2-C3}
  \begin{reviewercomment}
    국영문 요약에 제안한 방법에 대한 구체적인 내용을 기술하시기 바랍니다.
  \end{reviewercomment}

  \begin{authorresponse}
    국문 요약을 추가하고 영문 초록을 전면 보완했습니다. 두 요약 모두 (1) 기존
    CoC에서 빠질 수 있는 실행 허용 범위, (2) 일곱 요소의 자연어 execution guard,
    (3) 장면ㆍCoCㆍ가드ㆍ수치 후보를 입력받는 선택 절차, (4) 모델 외부 결정적
    검증, (5) guard만 바꾸는 paired-contract 평가를 구체적으로 기술합니다.
    핵심 결과로 표준 시간 평가의 위반률 0.319에서 0.000으로의 감소와 준수 목표
    완료율/paired accuracy 1.000, 미관측 시간값의 후보 정확도 0.635, 위반률
    0.351 및 paired accuracy 0.306을 함께 제시해 효과와 한계를 균형 있게
    요약했습니다.
  \end{authorresponse}

  \begin{manuscriptrevision}{국문 요약 및 영문 초록(1쪽)}
    국문 요약을 새로 추가하고 영문 초록을 동일한 구성으로 개정했습니다. 제안
    방법의 입력, 후보 선택, 모델 외부 검증과 쌍대 계약 평가를 구체화했습니다.
    핵심 수치와 ``실제 차량 안전 보장이 아님''이라는 주장 범위도 일관되게
    기술했습니다. 해당 개정은 보라색으로 표시했습니다.
  \end{manuscriptrevision}
\end{reviewitem}

\begin{reviewitem}{R2-C4}
  \begin{reviewercomment}
    학습된 가드 준수와 실제 실행안전의 구분을 한 이유가 무엇인지 구체적인 기술이
    필요합니다.
  \end{reviewercomment}

  \begin{authorresponse}
    두 개념을 구분한 이유는 평가 단위와 실패 원인이 다르기 때문입니다. 본 논문의
    learned guard compliance는 유한한 합성 후보와 연구자가 제공한 부분 계약에서
    목표와 가드를 만족하는 후보를 선택한 비율입니다. 실제 execution safety는
    센서 오인식, 계약에서 빠진 위험, 연속 궤적과 차량 동역학의 오차, 다른 도로
    사용자의 반응 및 지연된 상태 갱신까지 포함하는 폐루프 성질입니다. 따라서
    후보 준수율이 1.000이어도 가드가 불완전하거나 관측이 틀리면 실제 차량은
    위험할 수 있습니다. 이 구분은 합성 평가 점수를 안전 보장으로 확대하는 범주
    오류를 방지하고, 최신 상태를 사용하는 외부 검증기와 실행 차단장치가 별도로
    필요한 이유를 설명합니다.
  \end{authorresponse}

  \begin{manuscriptrevision}{제4.6절(7쪽), 제6.4절(13--14쪽)}
    학습된 가드 준수와 실제 실행 안전의 평가 단위 및 포함 실패 원인을 명시하고,
    두 개념을 구분해야 하는 이유를 추가했습니다. 폐루프 논의에서는 계약 상태의
    가역성과 최신성, 실행 차단의 역할을 이 구분과 연결했습니다. 해당 개정은
    보라색으로 표시했습니다.
  \end{manuscriptrevision}
\end{reviewitem}

\begin{reviewitem}{R2-C5}
  \begin{reviewercomment}
    연구 내용에 정확한 정보 전달을 위해서 문장 구성에 대한 점검도 필요합니다.
  \end{reviewercomment}

  \begin{authorresponse}
    원고 전체에서 용어와 문장 주체를 점검했습니다. ``기본 CoC'', ``제약 통합
    CoC'', ``실행 가드'', ``모델 외부 결정적 검증기'', ``가드 준수'' 및 ``실제
    실행 안전''을 일관되게 사용했습니다. 특히 연구 공백, 방법 입력ㆍ출력, 검증
    주체 및 결과의 해석 범위를 한 문장에 혼합하지 않도록 나누었습니다. 또한 각
    실험을 확인 목적, 비교 방법, 관찰 결과 및 의미로 분리하고, 서로 다른 모델의
    결과가 직접 비교되지 않는다는 주어와 범위를 명시했습니다.
  \end{authorresponse}

  \begin{manuscriptrevision}{제1--5절(1--7쪽), 제6.1절(13쪽)}
    대표적으로 서론의 연구 공백을 두 선행 연구 흐름과 본 연구의 역할로 분리하고,
    방법 절을 ``입력--대조 평가--외부 판정''의 세 질문으로 다시 설명했습니다.
    실험 절에는 증거의 순서를 안내하는 문단을 추가했으며, 논의에서는 관찰된
    수치와 허용 가능한 해석을 별도 문장으로 구분했습니다. 주요 보완 문장은
    보라색으로 표시했습니다.
  \end{manuscriptrevision}
\end{reviewitem}

\begin{reviewitem}{R2-C6}
  \begin{reviewercomment}
    제안하는 방법에 대해 또한 어떠한 장점이 있는지를 기술한다면 보다 완성도가
    높을 것 입니다.
  \end{reviewercomment}

  \begin{authorresponse}
    제안 방법의 장점을 네 가지로 명시했습니다. 첫째, 기존 CoC의 설명 구조를
    유지하면서 허용ㆍ금지ㆍ해제 조건을 자연어로 추가합니다. 둘째, paired-contract
    평가로 장면 암기와 guard에 따른 선택 전환을 구분합니다. 셋째, violation뿐
    아니라 goal completion, deadlock 및 excess conservatism을 함께 측정하여 항상
    정지하는 해법을 성공으로 보지 않습니다. 넷째, 모델의 후보 선택과 결정적 수치
    판정을 분리해 외부 verifier나 shield로 연결할 수 있습니다. 이러한 장점은
    합성 후보 선택의 해석 가능성과 진단 가능성에 관한 것이며 실제 차량의 안전
    보장으로 확대하지 않았습니다.
  \end{authorresponse}

  \begin{manuscriptrevision}{제1절(2쪽), 제6.1절(13쪽)}
    서론에 네 가지 장점을 명시하고, 논의에서 이 장점이 범용 주행 모델의 우월성이
    아니라 허용조건의 명시성, 반사실적 진단 및 독립 수치 판정에 있음을 기존
    CoC-only 조건과 대비했습니다. 해당 개정은 보라색으로 표시했습니다.
  \end{manuscriptrevision}
\end{reviewitem}

\begin{reviewitem}{R2-C7}
  \begin{reviewercomment}
    또한, 모델 외부 궤적 검증과 평가지표를 이용해서 제안하는 방법에 대한 성능이
    우수하다는 부분을 부각시킬 필요가 있습니다.
  \end{reviewercomment}

  \begin{authorresponse}
    외부 trajectory verifier는 모델의 자기 설명을 신뢰하지 않고 후보에 저장된
    정지 위치, 속도, 감속도, 가가속도, 간격 및 진입 시각을 guard 경계와 직접
    비교합니다. 이를 통해 candidate accuracy, violation, guard-compliant goal
    completion, deadlock, excess conservatism 및 paired-contract accuracy를 상호
    보완적으로 계산합니다. 표준 정적ㆍ시간ㆍ정지/차선 변경 평가에서 제약 통합
    CoC의 violation rate는 모두 0.000, paired-contract accuracy는 모두 1.000이었고
    목표 완료 관련 지표도 1.000이었습니다. 따라서 낮은 위반률이 동일 답 반복이나
    목표 포기의 결과가 아님을 강조했습니다. 동시에 미관측 시간값에서 위반률
    0.351, paired accuracy 0.306으로 악화된 결과를 함께 제시하여 우수성의 범위를
    표준 합성 평가로 한정했습니다.
  \end{authorresponse}

  \begin{manuscriptrevision}{제4.5절(6--7쪽), 제6.1절(13쪽), 제7절(15쪽)}
    방법 절에 외부 검증과 각 지표의 상호 보완적 진단 역할을 추가했습니다. 논의와
    결론에서는 표준 평가의 violation/goal completion/paired accuracy를 함께
    해석하여 제안 방법의 성능을 부각하고, 미관측 입력 결과를 근거로 적용 범위를
    제한했습니다. 해당 개정은 보라색으로 표시했습니다.
  \end{manuscriptrevision}
\end{reviewitem}

\section*{맺음말}

원고의 명확성과 엄밀성을 높일 수 있도록 세심한 의견을 주신 리뷰어 여러분께 다시
한번 감사드립니다. 개정 원고와 본 답변서가 제기된 모든 우려에 충실히 답하였기를
바랍니다.

감사합니다.\\
저자 일동

\end{document}
